\documentclass[11pt,a4paper]{article}

% --- Kodowanie i język ---
\usepackage[utf8]{inputenc}
\usepackage[T1]{fontenc}
\usepackage[polish]{babel}
\usepackage{lmodern}

% --- Układ strony ---
\usepackage[a4paper,margin=2.5cm]{geometry}
\usepackage{parskip}
\usepackage{microtype}

% --- Grafika, kolory, tabele ---
\usepackage{xcolor}
\usepackage{graphicx}
\usepackage{booktabs}
\usepackage{longtable}
\usepackage{array}
\usepackage{tabularx}
\usepackage{enumitem}

% --- Listingi kodu ---
\usepackage{listings}

% --- Hiperłącza ---
\usepackage[hidelinks]{hyperref}
\hypersetup{
    colorlinks=true,
    linkcolor=blue!50!black,
    urlcolor=blue!50!black,
    pdftitle={Dokumentacja techniczna -- FriendSync},
    pdfauthor={Kacper Skiba, Adam Młotkowski, Bartosz Sulecki}
}

% --- Definicje kolorów dla kodu ---
\definecolor{codebg}{rgb}{0.96,0.96,0.96}
\definecolor{codekeyword}{rgb}{0.13,0.29,0.53}
\definecolor{codecomment}{rgb}{0.40,0.55,0.40}
\definecolor{codestring}{rgb}{0.55,0.30,0.10}

\lstset{
    backgroundcolor=\color{codebg},
    basicstyle=\ttfamily\small,
    keywordstyle=\color{codekeyword}\bfseries,
    commentstyle=\color{codecomment}\itshape,
    stringstyle=\color{codestring},
    breaklines=true,
    showstringspaces=false,
    frame=single,
    rulecolor=\color{gray!40},
    tabsize=2,
    captionpos=b,
    columns=fullflexible,
    literate=%
        {ą}{{\k{a}}}1 {ć}{{\'c}}1 {ę}{{\k{e}}}1 {ł}{{\l{}}}1 {ń}{{\'n}}1
        {ó}{{\'o}}1 {ś}{{\'s}}1 {ż}{{\.z}}1 {ź}{{\'z}}1
        {Ą}{{\k{A}}}1 {Ć}{{\'C}}1 {Ę}{{\k{E}}}1 {Ł}{{\L{}}}1 {Ń}{{\'N}}1
        {Ó}{{\'O}}1 {Ś}{{\'S}}1 {Ż}{{\.Z}}1 {Ź}{{\'Z}}1
}

\newcolumntype{L}[1]{>{\raggedright\arraybackslash}p{#1}}

% --- Zrzuty ekranu: pokaz obrazek jesli istnieje, w przeciwnym razie placeholder ---
\graphicspath{{screeny/}{screenshots/}}
\newcommand{\screen}[2]{%
    \begin{figure}[h!]
        \centering
        \IfFileExists{screeny/#1}{%
            \includegraphics[width=0.85\textwidth,frame]{screeny/#1}%
        }{%
            \fbox{\parbox[c][4.5cm][c]{0.82\textwidth}{\centering\itshape
                Miejsce na zrzut ekranu --- dodaj plik \texttt{screeny/#1}}}%
        }
        \caption{#2}
    \end{figure}
}

% =====================================================================
\begin{document}

% --- Strona tytułowa ---
\begin{titlepage}
    \centering
    \vspace*{2cm}

    {\Huge\bfseries FriendSync\par}
    \vspace{0.5cm}
    {\Large Dokumentacja techniczna\par}

    \vspace{1.5cm}
    \rule{0.7\textwidth}{0.4pt}\par
    \vspace{0.5cm}
    {\large Aplikacja webowa do planowania spotkań ze znajomymi\\
    oraz rozliczania wspólnych wydatków grupowych\par}
    \vspace{0.5cm}
    \rule{0.7\textwidth}{0.4pt}\par

    \vfill

    {\large Autor: Kacper Skiba\par}
    \vspace{0.3cm}
    {\large \today\par}
\end{titlepage}

\tableofcontents
\newpage

% =====================================================================
\section{Opis projektu}

\textbf{FriendSync} to pełnostackowa aplikacja webowa wspierająca organizację wspólnych
wyjazdów i spotkań. Użytkownik może:

\begin{itemize}[noitemsep]
    \item zakładać konto i zarządzać profilem (avatar, bio, tagi),
    \item budować listę znajomych (zaproszenia, akceptacja, wyszukiwanie),
    \item tworzyć \emph{wydarzenia} (eventy) i zapraszać do nich znajomych,
    \item planować szczegóły wyjazdu poprzez \emph{podpunkty} (sub-events) z~harmonogramem,
    \item \emph{proponować lokalizacje} na mapie i \emph{głosować} na nie,
    \item prowadzić \emph{czat} wewnątrz wydarzenia,
    \item ewidencjonować \emph{wspólne wydatki} i automatycznie rozliczać długi
          z~minimalizacją liczby przelewów,
    \item otrzymywać \emph{powiadomienia} oraz widzieć \emph{status online} znajomych
          w~czasie rzeczywistym.
\end{itemize}

% =====================================================================
\section{Stos technologiczny}

\subsection{Backend}
\begin{longtable}{L{4cm} L{9cm}}
\toprule
\textbf{Technologia} & \textbf{Zastosowanie} \\
\midrule
\endhead
Python 3.12 & język programowania \\
FastAPI & framework webowy / REST API + WebSocket \\
Uvicorn & serwer ASGI \\
SQLAlchemy & ORM (mapowanie obiektowo-relacyjne) \\
PostgreSQL 16 & relacyjna baza danych \\
Pydantic v2 / pydantic-settings & walidacja danych i konfiguracja \\
PyJWT & tokeny JWT \\
passlib + bcrypt & hashowanie haseł \\
python-multipart & upload plików (avatary) \\
requests & integracja z Nominatim (geokodowanie) \\
\bottomrule
\end{longtable}

\subsection{Frontend}
\begin{longtable}{L{4cm} L{9cm}}
\toprule
\textbf{Technologia} & \textbf{Zastosowanie} \\
\midrule
\endhead
React 19 & biblioteka UI \\
Vite 8 & bundler / dev server \\
React Router 7 & routing po stronie klienta \\
Axios & klient HTTP \\
Tailwind CSS 4 & stylowanie \\
Leaflet + react-leaflet & mapy (OpenStreetMap) \\
date-fns / react-datepicker & obsługa dat \\
lucide-react & ikony \\
Vitest + Testing Library & testy jednostkowe \\
\bottomrule
\end{longtable}

\subsection{Infrastruktura}
\begin{itemize}[noitemsep]
    \item \textbf{Docker / Docker Compose} --- konteneryzacja,
    \item \textbf{Caddy} --- reverse proxy z~automatycznym HTTPS,
    \item \textbf{Nginx} --- serwowanie statycznego frontendu w~kontenerze.
\end{itemize}

% =====================================================================
\section{Architektura systemu}

Aplikacja ma architekturę \textbf{klient--serwer} z~wyraźnym rozdzieleniem na trójwarstwowy
backend (API $\rightarrow$ logika biznesowa $\rightarrow$ dane) oraz SPA po stronie klienta.

\begin{lstlisting}[language={}]
                   +---------------------------------------------+
                   |                  CADDY                      |
                   |   (reverse proxy, HTTPS, port 80/443)       |
                   +---------------+--------------+--------------+
                                   |              |
                   :80 (HTTPS)     |              |  :8443
              +--------------------v---+   +------v-------------------+
              |   FRONTEND (Nginx)     |   |   BACKEND (FastAPI)      |
              |   React SPA (build)    |   |   REST API + WebSocket   |
              +------------------------+   +-----------+-------------+
                                                       |
                                             +---------v----------+
                                             |   PostgreSQL 16    |
                                             +--------------------+
\end{lstlisting}

\textbf{Komunikacja:}
\begin{itemize}[noitemsep]
    \item \textbf{HTTP/REST} --- operacje CRUD (Axios $\rightarrow$ FastAPI),
    \item \textbf{WebSocket} (\texttt{/ws}) --- powiadomienia, status online, zdarzenia
          czasu rzeczywistego,
    \item \textbf{Nominatim (OpenStreetMap)} --- zewnętrzne API do zamiany adresu na
          współrzędne.
\end{itemize}

% =====================================================================
\section{Struktura repozytorium}

\begin{lstlisting}[language={}]
FriendSync/
|-- docker-compose.yml      # orkiestracja uslug (db, backend, frontend, caddy)
|-- Caddyfile               # konfiguracja reverse proxy + HTTPS
|-- README.md
|
|-- backend/
|   |-- Dockerfile
|   |-- requirements.txt
|   |-- pytest.ini
|   `-- app/
|       |-- main.py         # punkt wejscia, middleware, montaz routerow
|       |-- core/           # konfiguracja, baza, bezpieczenstwo, uploady
|       |-- models/         # modele ORM (SQLAlchemy)
|       |-- schemas/        # schematy Pydantic (walidacja I/O)
|       |-- crud/           # logika dostepu do danych / operacje biznesowe
|       |-- api/            # routery (endpointy REST + WebSocket)
|       |-- services/       # integracje zewnetrzne (geocoding)
|       `-- tests/          # testy pytest
|
`-- frontend/
    |-- Dockerfile
    |-- nginx.conf
    |-- vite.config.js
    `-- src/
        |-- App.jsx         # routing aplikacji
        |-- pages/          # widoki (AuthPage, Dashboard, EventDetails, ...)
        |-- components/     # komponenty wspoldzielone + Contexty
        `-- services/       # klient API, autoryzacja, preferencje
\end{lstlisting}

% =====================================================================
\section{Model danych}

System opiera się na relacyjnym modelu danych (katalog \texttt{backend/app/models/}).

\subsection{Diagram relacji (uproszczony)}

\begin{lstlisting}[language={}]
User  --<  Event (creator)
User  --<  EventParticipant  >--  Event
User  --<  Friendship  >--  User            (status: pending / accepted)
User  --<  EventInvitation  >--  Event      (status: pending / accepted / declined)
Event --<  SubEvent  >--  LocationProposal  (opcjonalnie)
Event --<  LocationProposal  --<  LocationVote  >--  User
Event --<  Message  >--  User (author)
Event --<  Expense (payer=User)  --<  ExpenseShare  >--  User
User  --<  Notification
\end{lstlisting}

\subsection{Opis encji}

\begin{longtable}{L{2.6cm} L{2.8cm} L{7.6cm}}
\toprule
\textbf{Encja} & \textbf{Tabela} & \textbf{Kluczowe pola} \\
\midrule
\endhead
User & \texttt{users} & \texttt{username} (unikalne), \texttt{email} (unikalne),
    \texttt{password\_hash}, \texttt{profile\_image}, \texttt{bio}, \texttt{tags},
    \texttt{last\_active} \\
Event & \texttt{events} & \texttt{title}, \texttt{description}, \texttt{creator\_id},
    \texttt{event\_date} \\
EventParticipant & \texttt{event\_participants} & \texttt{event\_id}, \texttt{user\_id},
    \texttt{role} (admin/member), \texttt{joined\_at} \\
EventInvitation & \texttt{event\_invitations} & \texttt{event\_id}, \texttt{user\_id}
    (zapraszany), \texttt{inviter\_id}, \texttt{status} \\
SubEvent & \texttt{sub\_events} & \texttt{event\_id}, \texttt{title}, \texttt{start\_time},
    \texttt{location\_id} (opcjonalne podpięcie punktu na mapie) \\
LocationProposal & \texttt{location\_proposals} & \texttt{event\_id}, \texttt{creator\_id},
    \texttt{name}, \texttt{latitude}, \texttt{longitude} \\
LocationVote & \texttt{location\_votes} & \texttt{location\_id}, \texttt{user\_id},
    \texttt{vote\_value} \\
Message & \texttt{messages} & \texttt{event\_id}, \texttt{author\_id}, \texttt{content} \\
Expense & \texttt{expenses} & \texttt{event\_id}, \texttt{payer\_id}, \texttt{amount},
    \texttt{description} \\
ExpenseShare & \texttt{expense\_shares} & \texttt{expense\_id}, \texttt{user\_id},
    \texttt{amount}, \texttt{is\_settled}, \texttt{settled\_at} \\
Friendship & \texttt{friendships} & \texttt{user\_id}, \texttt{friend\_id},
    \texttt{status} \\
Notification & \texttt{notifications} & \texttt{user\_id}, \texttt{type}, \texttt{message},
    \texttt{is\_read} \\
\bottomrule
\end{longtable}

\subsection{Kaskady i integralność}
\begin{itemize}[noitemsep]
    \item usunięcie \texttt{Event} kaskadowo usuwa uczestników, wydatki, lokalizacje,
          wiadomości i~podpunkty (\texttt{cascade="all, delete-orphan"}),
    \item usunięcie \texttt{LocationProposal} usuwa powiązane głosy,
    \item usunięcie \texttt{LocationProposal} powiązanej z~\texttt{SubEvent} ustawia
          \texttt{location\_id = NULL} (\texttt{ON DELETE SET NULL}) --- podpunkt nie znika,
          traci tylko powiązanie z~mapą.
\end{itemize}

\textbf{Migracje:} projekt nie używa Alembica. Drobne, idempotentne migracje
(\texttt{ADD COLUMN IF NOT EXISTS}) wykonywane są automatycznie przy starcie aplikacji
w~\texttt{main.py} (funkcja \texttt{\_run\_lightweight\_migrations()}). Schemat tabel tworzony
jest przez \texttt{Base.metadata.create\_all(bind=engine)}.

% =====================================================================
\section{Backend --- warstwy i logika}

Backend stosuje rozdzielenie odpowiedzialności na cztery warstwy:

\begin{enumerate}[noitemsep]
    \item \textbf{\texttt{api/} (routery)} --- definicja endpointów, walidacja wejścia,
          autoryzacja (zależność \texttt{get\_current\_user}), wywoływanie warstwy CRUD,
          serializacja odpowiedzi.
    \item \textbf{\texttt{schemas/} (Pydantic)} --- kontrakty danych wejściowych/wyjściowych,
          walidacja typów i~ograniczeń (np. \texttt{amount > 0}).
    \item \textbf{\texttt{crud/} (logika biznesowa)} --- operacje na bazie, reguły domenowe
          (walidacja sum w~wydatkach, algorytm rozliczeń).
    \item \textbf{\texttt{models/} (ORM)} --- mapowanie obiektowo-relacyjne SQLAlchemy.
\end{enumerate}

\textbf{Warstwy pomocnicze:}
\begin{itemize}[noitemsep]
    \item \texttt{core/config.py} --- konfiguracja przez zmienne środowiskowe
          (pydantic-settings), walidacja \texttt{SECRET\_KEY},
    \item \texttt{core/database.py} --- silnik SQLAlchemy, fabryka sesji, zależność
          \texttt{get\_db()},
    \item \texttt{core/security.py} --- hashowanie haseł, polityka haseł, generowanie
          tokenów JWT,
    \item \texttt{core/uploads.py} --- bezpieczny zapis avatarów,
    \item \texttt{services/geocoding.py} --- geokodowanie adresów przez Nominatim.
\end{itemize}

% =====================================================================
\section{API REST --- wykaz endpointów}

Większość endpointów wymaga nagłówka \texttt{Authorization: Bearer <token>}.

\subsection{Użytkownicy --- \texttt{/api/users}}
\begin{longtable}{L{1.8cm} L{5cm} L{6cm}}
\toprule
\textbf{Metoda} & \textbf{Ścieżka} & \textbf{Opis} \\
\midrule
\endhead
POST & \texttt{/register} & rejestracja (multipart: email, username, password, avatar) \\
POST & \texttt{/login} & logowanie (OAuth2), zwraca token JWT \\
GET & \texttt{/me} & dane zalogowanego użytkownika \\
PATCH & \texttt{/me} & edycja profilu (zmiana e-maila/hasła wymaga obecnego hasła) \\
GET & \texttt{/me/finances/summary} & globalne podsumowanie długów i~należności \\
DELETE & \texttt{/me} & trwałe usunięcie konta wraz z~danymi \\
GET & \texttt{/\{user\_id\}/online-status} & status online użytkownika \\
\bottomrule
\end{longtable}

\subsection{Wydarzenia --- \texttt{/api/events}}
\begin{longtable}{L{1.8cm} L{5.5cm} L{5.5cm}}
\toprule
\textbf{Metoda} & \textbf{Ścieżka} & \textbf{Opis} \\
\midrule
\endhead
POST & \texttt{""} & utworzenie wydarzenia \\
GET & \texttt{""} & lista wydarzeń użytkownika \\
GET & \texttt{/\{id\}/participants} & uczestnicy \\
POST & \texttt{/\{id\}/invite} & zaproszenie znajomego \\
GET & \texttt{/invitations} & otrzymane zaproszenia \\
POST & \texttt{/invitations/\{id\}/accept|decline} & akceptacja / odrzucenie \\
POST/GET & \texttt{/\{id\}/messages} & czat (wysłanie / lista) \\
POST/GET & \texttt{/\{id\}/expenses} & wydatki (dodanie / lista) \\
POST & \texttt{/\{id\}/shares/\{sid\}/settle} & spłata pojedynczego udziału \\
POST & \texttt{/\{id\}/creditors/\{cid\}/settle-all} & zbiorcza spłata wobec wierzyciela \\
GET & \texttt{/\{id\}/finances/summary} & bilans i~lista przelewów \\
POST/PUT/DELETE & \texttt{/\{id\}/sub-events} & zarządzanie podpunktami \\
PUT/DELETE & \texttt{/\{id\}} & edycja / usunięcie wydarzenia \\
DELETE & \texttt{/\{id\}/participants/\{uid\}} & usunięcie uczestnika \\
PUT & \texttt{/\{id\}/transfer-ownership/\{nid\}} & przekazanie własności \\
\bottomrule
\end{longtable}

\subsection{Znajomi --- \texttt{/api/friends}}
\begin{longtable}{L{1.8cm} L{5cm} L{6cm}}
\toprule
\textbf{Metoda} & \textbf{Ścieżka} & \textbf{Opis} \\
\midrule
\endhead
POST & \texttt{/request} & wysłanie zaproszenia \\
POST & \texttt{/\{id\}/accept} & akceptacja \\
GET & \texttt{""} & lista znajomych \\
GET & \texttt{/pending} & oczekujące zaproszenia \\
GET & \texttt{/search-users} & wyszukiwanie użytkowników \\
DELETE & \texttt{/\{friend\_id\}} & usunięcie znajomego \\
\bottomrule
\end{longtable}

\subsection{Lokalizacje i powiadomienia}
\begin{longtable}{L{1.8cm} L{6cm} L{5cm}}
\toprule
\textbf{Metoda} & \textbf{Ścieżka} & \textbf{Opis} \\
\midrule
\endhead
GET/POST & \texttt{/api/events/\{id\}/locations} & propozycje lokalizacji \\
DELETE & \texttt{/api/locations/\{id\}} & usunięcie propozycji \\
POST/DELETE & \texttt{/api/locations/\{id\}/votes} & głos / wycofanie głosu \\
GET & \texttt{/api/notifications} & lista powiadomień \\
PUT & \texttt{/api/notifications/\{id\}/read} & oznaczenie jako przeczytane \\
DELETE & \texttt{/api/notifications} & wyczyszczenie powiadomień \\
\bottomrule
\end{longtable}

\noindent Pełna, interaktywna specyfikacja (OpenAPI/Swagger) dostępna jest automatycznie
pod adresem \texttt{/docs} dzięki FastAPI.

% =====================================================================
\section{Komunikacja w czasie rzeczywistym (WebSocket)}

Endpoint \texttt{/ws?token=<JWT>} (\texttt{api/websocket.py}) obsługuje połączenia czasu
rzeczywistego.

\begin{itemize}[noitemsep]
    \item \textbf{Autoryzacja:} token JWT przekazywany w~query stringu; brak lub nieważny
          token skutkuje zamknięciem połączenia kodem \texttt{1008}.
    \item \textbf{\texttt{ConnectionManager}} utrzymuje mapę
          \texttt{user\_id $\rightarrow$ [WebSocket]} (jeden użytkownik może mieć wiele
          kart/urządzeń).
    \item \textbf{Status online} jest wyznaczany \emph{wyłącznie} na podstawie aktywnego
          połączenia WebSocket (\texttt{is\_user\_online}), a~nie pola \texttt{last\_active}
          --- świadoma decyzja projektowa (poprzednie podejście resetowało zegar przy każdym
          żądaniu i~użytkownik nigdy nie przechodził w~offline).
    \item Przy połączeniu/rozłączeniu manager powiadamia znajomych o~zmianie statusu.
\end{itemize}

\textbf{Typy komunikatów} rozsyłanych do klientów:
\begin{itemize}[noitemsep]
    \item \texttt{user\_status} --- zmiana statusu online znajomego,
    \item \texttt{profile\_updated} --- znajomy zaktualizował profil,
    \item \texttt{friend\_removed} --- znajomy usunął konto / znajomość,
    \item powiadomienia o~zdarzeniach w~wydarzeniach (\texttt{broadcast\_to\_event}).
\end{itemize}

% =====================================================================
\section{Moduł finansowy --- algorytm rozliczeń}

Najbardziej złożona część logiki biznesowej (\texttt{crud/expense.py}).

\subsection{Model wydatku}
Każdy wydatek (\texttt{Expense}) ma płatnika (\texttt{payer\_id}) i~listę udziałów
(\texttt{ExpenseShare}) --- kto i~ile jest winny. Walidacja przy tworzeniu wymaga m.in.:
\begin{itemize}[noitemsep]
    \item kwoty dodatniej,
    \item sumy udziałów równej kwocie wydatku (z~tolerancją jednego grosza),
    \item braku duplikatów użytkowników w~udziałach,
    \item przynależności płatnika i~dłużników do wydarzenia.
\end{itemize}
Jeśli płatnik występuje też we własnych udziałach, jego udział jest od razu oznaczany jako
spłacony (\texttt{is\_settled=True}) --- nie płaci samemu sobie.

\subsection{Obliczanie bilansu}
Funkcja \texttt{\_compute\_event\_balances} liczy saldo per użytkownik z~\textbf{niespłaconych}
udziałów: płatnik dostaje $+amount$, dłużnik $-amount$. Wpisy audytowe spłat (rozpoznawane po
prefiksie opisu \texttt{"Rozliczenie"}) są pomijane, by nie zaburzać salda.

\subsection{Minimalizacja liczby przelewów (greedy min-cash-flow)}
Funkcja \texttt{calculate\_finance\_summary} zamienia bilans na \textbf{minimalną liczbę
przelewów}:
\begin{enumerate}[noitemsep]
    \item dzieli użytkowników na dłużników (saldo ujemne) i~wierzycieli (saldo dodatnie),
    \item sortuje obie listy malejąco po kwocie,
    \item dopasowuje największego dłużnika do największego wierzyciela, generując przelew na
          $\min(\text{dług}, \text{należność})$, i~powtarza aż do wyrównania.
\end{enumerate}
Dzięki temu zamiast „każdy z~każdym” powstaje krótka lista konkretnych przelewów.

\subsection{Spłaty}
\begin{itemize}[noitemsep]
    \item \texttt{settle\_share} --- oznacza pojedynczy udział jako spłacony i~tworzy
          \textbf{wpis audytowy} (oddzielny \texttt{Expense} z~prefiksem \texttt{"Rozliczenie"})
          jako ślad w~historii.
    \item \texttt{settle\_all\_with\_creditor} --- zbiorczo zamyka wszystkie otwarte długi
          użytkownika wobec jednego wierzyciela.
    \item \texttt{calculate\_global\_finance\_summary} --- agreguje długi/należności
          użytkownika ze wszystkich wydarzeń (widok globalny na Dashboardzie).
\end{itemize}

\noindent Operacje na pieniądzach używają stałej tolerancji \texttt{CENT = 0.01} z~uwagi na
niedokładność arytmetyki zmiennoprzecinkowej.

% =====================================================================
\section{Bezpieczeństwo}

\subsection{Uwierzytelnianie i hasła}
\begin{itemize}[noitemsep]
    \item hasła hashowane algorytmem \textbf{bcrypt} (passlib),
    \item polityka haseł: 10--128 znaków, min. 3 z~4 klas znaków, blokada haseł pospolitych
          (\texttt{validate\_password\_strength}),
    \item \textbf{JWT} (HS256) z~polami \texttt{exp}, \texttt{iat}, \texttt{jti}, \texttt{type}
          --- krótki czas życia (60 min),
    \item \textbf{ochrona przed user enumeration}: przy logowaniu nieistniejącego konta i~tak
          wykonywany jest pełny koszt bcrypt (\texttt{dummy\_verify}) --- wyrównanie czasu
          odpowiedzi (timing attack),
    \item zmiana e-maila/hasła wymaga potwierdzenia obecnym hasłem (re-auth).
\end{itemize}

\subsection{Konfiguracja}
W trybie \texttt{production} aplikacja \textbf{nie wystartuje} bez silnego \texttt{SECRET\_KEY}
($\geq$ 32 znaki); w~\texttt{development} generowany jest losowy klucz efemeryczny.

\subsection{Nagłówki bezpieczeństwa (SecurityHeadersMiddleware)}
\begin{itemize}[noitemsep]
    \item \texttt{X-Content-Type-Options: nosniff},
    \item \texttt{X-Frame-Options: DENY},
    \item \texttt{Referrer-Policy: strict-origin-when-cross-origin},
    \item \texttt{Permissions-Policy} (ograniczenie geolokalizacji/kamery/mikrofonu),
    \item \texttt{Strict-Transport-Security} (HSTS) --- tylko w~produkcji.
\end{itemize}

\subsection{Pozostałe}
\begin{itemize}[noitemsep]
    \item \textbf{CORS} --- restrykcyjna allowlista origin (\texttt{ALLOWED\_ORIGINS}),
          ograniczone metody i~nagłówki,
    \item \textbf{Upload plików} --- walidacja rozmiaru avatara (limit 5 MB,
          \texttt{MAX\_AVATAR\_SIZE}).
\end{itemize}

% =====================================================================
\section{Frontend}

SPA w~React 19 z~routingiem klienckim (\texttt{App.jsx}):

\begin{longtable}{L{4cm} L{3.5cm} L{5.5cm}}
\toprule
\textbf{Ścieżka} & \textbf{Widok} & \textbf{Rola} \\
\midrule
\endhead
\texttt{/} & AuthPage & logowanie / rejestracja \\
\texttt{/dashboard} & Dashboard & lista wydarzeń, znajomi, mapa globalna, finanse \\
\texttt{/events/:id} & EventDetails & szczegóły, czat, mapa, podpunkty \\
\texttt{/events/:id/finance} & EventFinance & wydatki i~rozliczenia \\
\texttt{/edit-profile} & EditProfilePage & edycja profilu \\
\texttt{/settings} & SettingsPage & ustawienia (waluta, wygląd) \\
\bottomrule
\end{longtable}

\textbf{Zarządzanie stanem globalnym} --- React Context:
\begin{itemize}[noitemsep]
    \item \texttt{WebSocketContext} --- pojedyncze, współdzielone połączenie WebSocket,
    \item \texttt{CurrencyContext} --- wybrana waluta,
    \item \texttt{DialogContext} --- globalne okna dialogowe / potwierdzenia.
\end{itemize}

\textbf{Warstwa usług} (\texttt{src/services/}):
\begin{itemize}[noitemsep]
    \item \texttt{api.js} --- skonfigurowany klient Axios (bazowy URL, dołączanie tokenu),
    \item \texttt{authService.js} --- logowanie/wylogowanie, obsługa tokenu,
    \item \texttt{preferences.js} --- preferencje wyglądu zapisywane lokalnie.
\end{itemize}

\textbf{Mapy} --- Leaflet + react-leaflet (\texttt{EventMapComponent},
\texttt{GlobalDashboardMap}) na danych OpenStreetMap; współrzędne propozycji uzyskiwane
z~geokodowania po stronie backendu.

\subsection{Zgodność z~heurystykami Nielsena}

Interfejs FriendSync był projektowany w~oparciu o~10 heurystyk użyteczności Jakoba Nielsena.
Poniżej zestawienie każdej heurystyki z~jej realizacją w~aplikacji:

\begin{longtable}{L{4.5cm} L{8.5cm}}
\toprule
\textbf{Heurystyka} & \textbf{Realizacja w~FriendSync} \\
\midrule
\endhead
1. Widoczność stanu systemu & status online znajomych w~czasie rzeczywistym (WebSocket),
    wskaźniki ładowania, liczniki nieprzeczytanych powiadomień, potwierdzenia akcji (toast). \\
2. Zgodność z~rzeczywistością & język naturalny i~dziedzinowy (znajomi, wydarzenia,
    wydatki, „kto komu jest winny”), metafory mapy i~kalendarza zamiast pojęć technicznych. \\
3. Kontrola i~swoboda użytkownika & możliwość anulowania/odrzucenia zaproszeń, wycofania
    głosu na lokalizację, edycji i~usuwania wydarzeń, opuszczenia wydarzenia. \\
4. Spójność i~standardy & jednolite komponenty (\texttt{DialogContext}), spójna paleta
    Tailwind, powtarzalny układ widoków i~ikonografia (\texttt{lucide-react}). \\
5. Zapobieganie błędom & walidacja formularzy po stronie klienta i~serwera (Pydantic),
    okna potwierdzenia dla akcji destrukcyjnych, polityka silnych haseł. \\
6. Rozpoznawanie zamiast przypominania & widoczne listy wydarzeń, znajomych i~zaproszeń,
    podpowiedzi w~wyszukiwaniu użytkowników, mapa z~oznaczonymi propozycjami. \\
7. Elastyczność i~wydajność & skróty: zbiorcza spłata wobec wierzyciela
    (\texttt{settle-all}), globalny widok finansów na Dashboardzie, zapisywane preferencje. \\
8. Estetyka i~minimalizm & czysty, responsywny layout, prezentacja tylko istotnych
    informacji, podział finansów na widok per wydarzenie i~globalny. \\
9. Pomoc w~rozpoznaniu i~naprawie błędów & czytelne komunikaty błędów (np. niepoprawna suma
    udziałów, błędne hasło) zamiast kodów technicznych. \\
10. Pomoc i~dokumentacja & opisowe etykiety i~podpowiedzi, intuicyjny przepływ
    (rejestracja $\rightarrow$ znajomi $\rightarrow$ wydarzenie $\rightarrow$ rozliczenie),
    interaktywna dokumentacja API pod \texttt{/docs}. \\
\bottomrule
\end{longtable}

% =====================================================================
\section{Zrzuty ekranu}

Poniżej zaprezentowano kluczowe widoki interfejsu użytkownika aplikacji FriendSync.

\medskip
\noindent\textit{Uwaga: pliki graficzne należy umieścić w~katalogu \texttt{screeny/} obok pliku
\texttt{.tex}, pod nazwami podanymi w~ramkach. Do czasu ich dodania w~dokumencie wyświetlane
są miejsca zastępcze.}

\screen{auth.png}{Ekran logowania i~rejestracji (\texttt{AuthPage}).}

\screen{dashboard.png}{Pulpit użytkownika --- lista wydarzeń, znajomi, mapa globalna i~finanse.}

\screen{event-details.png}{Szczegóły wydarzenia --- czat, podpunkty i~mapa propozycji lokalizacji.}

\screen{event-map.png}{Propozycje lokalizacji na mapie wraz z~głosowaniem.}

\screen{event-finance.png}{Moduł finansowy --- wydatki i~rozliczenia z~minimalizacją przelewów.}

\screen{edit-profile.png}{Edycja profilu użytkownika (avatar, bio, tagi).}

\screen{settings.png}{Ustawienia aplikacji (waluta, wygląd).}

% =====================================================================
\section{Testy}

\subsection{Backend (pytest)}
\begin{itemize}[noitemsep]
    \item \texttt{conftest.py} --- fixtury, izolowana baza testowa,
    \item \texttt{test\_api.py} --- testy integracyjne endpointów,
    \item \texttt{test\_finance.py} --- testy algorytmu rozliczeń (bilans, minimalizacja
          przelewów, spłaty),
    \item \texttt{test\_schemas.py} --- walidacja schematów Pydantic,
    \item \texttt{test\_security.py} --- hashowanie haseł, polityka haseł, JWT.
\end{itemize}

\begin{lstlisting}[language=bash]
cd backend
pytest
\end{lstlisting}

\subsection{Frontend (Vitest + Testing Library)}
Testy logiki usług: \texttt{authService.test.js}, \texttt{preferences.test.js}.

\begin{lstlisting}[language=bash]
cd frontend
npm run test
\end{lstlisting}

% =====================================================================
\section{Wdrożenie (Docker / Caddy)}

Plik \texttt{docker-compose.yml} definiuje cztery usługi:

\begin{longtable}{L{2.2cm} L{4cm} L{6.5cm}}
\toprule
\textbf{Usługa} & \textbf{Obraz / źródło} & \textbf{Rola} \\
\midrule
\endhead
\texttt{db} & \texttt{postgres:16} & baza danych (wolumen \texttt{postgres\_data}) \\
\texttt{backend} & build \texttt{./backend} & FastAPI + Uvicorn (port 8000) \\
\texttt{frontend} & build \texttt{./frontend} & React zbudowany i~serwowany przez Nginx \\
\texttt{caddy} & \texttt{caddy:alpine} & reverse proxy + HTTPS (porty 80/443/8443) \\
\bottomrule
\end{longtable}

\textbf{Caddy} (\texttt{Caddyfile}) kieruje:
\begin{itemize}[noitemsep]
    \item \texttt{friendsync.me} $\rightarrow$ frontend (Nginx),
    \item \texttt{friendsync.me:8443} $\rightarrow$ backend (API + WebSocket),
    \item automatycznie zarządza certyfikatami TLS.
\end{itemize}

\textbf{Backend Dockerfile} --- \texttt{python:3.12-slim}, instalacja zależności systemowych
(\texttt{gcc}, \texttt{libpq-dev}) i~pythonowych.

\textbf{Frontend Dockerfile} --- wieloetapowy build: \texttt{node:22-alpine} buduje aplikację
Vite, następnie artefakty kopiowane są do obrazu \texttt{nginx:alpine}.

\textbf{Zmienne środowiskowe} (przez \texttt{.env}): \texttt{DB\_PASSWORD},
\texttt{SECRET\_KEY}, \texttt{ALLOWED\_ORIGINS}, \texttt{VITE\_API\_URL},
\texttt{DATABASE\_URL}.

\begin{lstlisting}[language=bash]
docker compose up --build
\end{lstlisting}

% =====================================================================
\section{Uruchomienie lokalne}

\subsection{Backend}
\begin{lstlisting}[language=bash]
cd backend
python -m venv venv
# Windows:
venv\Scripts\activate
# Linux/Mac:
source venv/bin/activate

pip install -r requirements.txt
uvicorn app.main:app --reload
\end{lstlisting}
API dostępne pod \texttt{http://localhost:8000}, dokumentacja Swagger pod
\texttt{http://localhost:8000/docs}.

\subsection{Frontend}
\begin{lstlisting}[language=bash]
cd frontend
npm install
npm run dev
\end{lstlisting}
Aplikacja dostępna pod \texttt{http://localhost:5173}.

\vspace{0.5cm}
\noindent\textit{W trybie deweloperskim, jeśli \texttt{SECRET\_KEY} nie jest ustawiony, backend
generuje losowy klucz przy każdym starcie --- wymaga to ponownego logowania po restarcie
serwera.}

\end{document}
